% Stand 21. Juli 2026
Die zentralen, im Zuge dieser Arbeit zu adressierenden Fragen
lauten:
\begin{itemize}[topsep=2pt,parsep=1pt,itemsep=0.5pt]
\item[a) ] liefern die betrachteten Feinstaubdaten für den Großraum München 
Hinweise auf eine mögliche Anisotropie und
\item[b) ] können die Feinstaubdaten unter Nutzung eines entsprechenden
Kovarianz-Modells mithilfe des Kriging hinreichend gut 
beschrieben werden.
\end{itemize}

\vspace*{6pt}
Es ist wichtig, sich zu vergegenwärtigen, dass die Frage nach der
Gültigkeit der (in ihren mathematischen Konsequenzen angenehmeren) 
Isotropie-Annahme zu klären ist, \textit{bevor} man die Methode 
des Kriging anwendet: die Wahl des am besten geeigneten
Kovarianz-Modells für den Kriging Algorithmus beinhaltet 
neben der Modell-Funktion auch die Vorentscheidung hinsichtlich 
einer isotropen oder anisotropen Struktur der Daten, sowie 
ggf.\,die Modellierung der Anisotropie 
(vgl.\,Kapitel \ref{KapitelTheorieAnisotropie}). 

Auf diesem Hintergrund werden nachfolgend 
\begin{itemize}[topsep=2pt,parsep=1pt,itemsep=0.5pt]
\item[$\circ$ ] im Kapitel 
\ref{KapitelDatenauswertung_Ergebnisse_Beobachtungen}
zunächst allgemeine Beobachtung bzgl.\,räumlicher und 
statistischer Verteilung der Messdaten herausgearbeitet.
\item[$\circ$ ] Im Kapitel
\ref{KapitelDatenauswertung_Ergebnisse_Variogramme}
werden die Daten mittels 
direktionaler Variogramme analysiert,
welche die Standardmethode der Geostatistik zur 
Aufdeckung einer möglichen anisotropen Datenstruktur
darstellen.
\item[$\circ$ ] Anschließend wird in 
Kapitel \ref{KapitelDatenauswertung_Ergebnisse_Kriging} 
die Frage nach einer möglichen 
Anisotropie mithilfe von 
konstruierten anisotropen Kovarianz-Modellen unter Nutzung 
des Kriging Interpolationsverfahren beleuchtet bevor 
\item[$\circ$] abschließend in Kapitel 
\ref{KapitelDatenauswertung_VerwandteArbeiten} die 
gewonnenen Einsichten in den 
Kontext ähnlicher Arbeiten zum Thema der räumlichen 
Analyse von Feinstaubmessungen gestellt werden.
\end{itemize} 

\subsection{Ergebnisse}
\label{KapitelDatenauswertung_Ergebnisse}
Alle nachfolgend gezeigten Auswertungen und Grafiken
wurden in der Programmiersprache \texttt{Python} erstellt. 
Hierbei wurden die Analysen zu ungerichteten und direktionalen
Variogrammen mithilfe der Bibliothek \texttt{Scikit GStat}
($\rightsquigarrow$ \href{https://scikit-gstat.readthedocs.io/en/latest/index.html}{Scikit GStat Home}) und Analysen mit 
Bezug zu Kovarianz-Modellen und Kriging mithilfe der 
Bibliothek \texttt{GSTools} 
($\rightsquigarrow$ 
\href{https://geostat-framework.readthedocs.io/projects/gstools/en/stable/index.html}{GSTools Quickstart}) vorgenommen.\\
%Es wurde ausschließlich die Variante des ordinary Kriging
%(vgl.\,Kapitel \ref{KapitelTheorie_OrdinaryKriging})
%für die Auswertungen herangezogen.

Insgesamt standen für die Auswertung die täglichen, 
bereits bereinigten Daten und an verschiedenen Orten 
positionierter Sensoren für 
die Monate Juli 2025 (96 Sensoren), November 2025 (87 Sensoren)
 und April 2026 (83 Sensoren) zur
Verfügung.\\
Die erratisch anmutende Auswahl der betrachteten Monate rührt daher, dass 
anfänglich auch der \textit{saisonale Einfluss} auf die Modellierung der
Feinstaubdaten untersucht und daher ein
exemplarischer Vertreter aller vier Jahreszeiten einbezogen 
werden sollte. Aus Zeitgründen wurde kein Winter-Monat mehr 
berücksichtigt und diese Untersuchungen als 
Perspektive in Kapitel \ref{KapitelDatenauswertung_Diskussion} aufgenommen.  \\   

Den typischen Verlauf der $\mathrm{P}_{2.5}$-Daten
eines einzelnen Sensors am 1.11.2025
zeigt Abb.\,(\ref{Abb_Datenauswertung_Sensor4430_P25TagesMessungNov2025}); 
der Verlauf desselben Sensors für den gesamten 
November ist in 
Abb.\,(\ref{Abb_Datenauswertung_Sensor4430_P25MonatsMessungNov2025})
dargestellt. 

\begin{figure}[H] 
\centering
\includegraphics[scale = 0.45]{Datenauswertung_Sensor4430_P25TagesMessungNov2025} 
\caption{\small{Tagesverlauf der $\mathrm{P}_{2.5}$-Messungen des 
Sensors (ID 4430) am 1.\,November 2025; 
eingezeichnet der Mittelwert (\textcolor{NavyBlue}{\rule[1mm]{3mm}{0.5mm}})
und der Median (\textcolor{CarnationPink}
{\rule[1mm]{3mm}{0.5mm}})}}
\label{Abb_Datenauswertung_Sensor4430_P25TagesMessungNov2025}
\end{figure} 

\begin{figure}[H] 
\centering
\includegraphics[scale = 0.45]{Datenauswertung_Sensor4430_P25MonatsMessungNov2025} 
\caption{\small{Monatsverlauf der $\mathrm{P}_{2.5}$-Messungen des 
Sensors (ID 4430) im November 2025; 
eingezeichnet der Mittelwert (\textcolor{NavyBlue}{\rule[1mm]{3mm}{0.5mm}})
und der Median (\textcolor{CarnationPink}
{\rule[1mm]{3mm}{0.5mm}})}}
\label{Abb_Datenauswertung_Sensor4430_P25MonatsMessungNov2025}
\end{figure} 

Da in beiden Abbildungen eine hohe Schwankungsbreite
erkennbar ist, wurden für die nachfolgenden Betrachtungen die
monatlichen arithmetischen Mittelwerte und der monatliche Median 
pro Sensor betrachtet\footnote{Ein typischer 
Sensor nimmt pro Tag etwa 580 Messungen 
(jeweils für $\mathrm{P}_{10}$ und $\mathrm{P}_{2.5}$ ) auf, so dass 
für die Berechnung des Mittelwerts bzw.\,des Medians 
rund $30\cdot 580=17400$ Werte herangezogen werden.}.

\subsubsection{Allgemeine Beobachtungen}
\label{KapitelDatenauswertung_Ergebnisse_Beobachtungen}
Abb.\,(\ref{Abb_Datenauswertung_VerteilungMesspunkteApr2026}) 
zeigt beispielhaft für den April 2026 die Verteilung der 83 
Sensoren innerhalb der Stadtgrenzen von München. Hierbei 
symbolisieren die braun eingezeichneten Linien primäre 
und sekundäre Straßen gemäß der
\texttt{OpenStreetMap} Einteilung
($\rightsquigarrow$ \href{https://www.openstreetmap.org/relation/62428}
{OSM: Munich}).

\begin{figure}[H] 
\centering
\includegraphics[scale = 0.52]{Datenauswertung_VerteilungMesspunkteApr2026} 
\caption{\small{Räumliche Verteilung der für die Analyse der $\mathrm{P}_{2.5}$-Messungen 
im April 2026 berücksichtigten Sensoren (\textcolor{ForestGreen}{$\bullet$}); Zentrum der
azimuthalen Projektion ist der Marienplatz 
(\textcolor{gray}{$\boldsymbol{\mathrm{X}}$})}}
\label{Abb_Datenauswertung_VerteilungMesspunkteApr2026}
\end{figure} 
 
\begin{figure}[H] 
\centering
\includegraphics[scale = 0.52]{Datenauswertung_P25MessungenApr2026} 
\caption{\small{Arithmetischer Mittelwert der $\mathrm{P}_{2.5}$-Messungen 
im April 2026; Zentrum der azimuthalen Projektion ist der 
Marienplatz (\textcolor{gray}{$\boldsymbol{\mathrm{X}}$})}}
\label{Abb_Datenauswertung_P25MessungenApr2026}
\end{figure} 

In 
Abb.\,(\ref{Abb_Datenauswertung_P25MessungenApr2026}) 
sind farbkodiert die arithmetischen Mittelwerte der 
$\mathrm{P}_{2.5}$-Messungen dieser Sensoren wiedergegeben.
Bei der Betrachtung erkennt man, dass die Sensoren auf
 einer Fläche von 
überschlagen etwa $20\,km$ Ausdehnung in Ost-West-Richtung und 
$15\,km$ 
Ausdehnung in Nord-Süd-Richtung verteilt sind.
Daraus ergibt sich 
die Abschätzung, dass die Verteilung der paarweisen Distanzen nach 
oben durch einen Wert von etwa $\sqrt{20^2 + 15^2}= 25\,km$ beschränkt
sein sollte. 

\begin{figure}[H] 
\centering
\includegraphics[scale = 0.45]{Datenauswertung_VerteilungDistanzenApr2026} 
\caption{\small{Häufigkeits-Verteilung der paarweisen Distanzen 
$R$ der 
83 Sensoren im April 2026}}
\label{Abb_Datenauswertung_VerteilungDistanzenApr2026}
\end{figure} 

Diese Abschätzung wird durch die 
Abb.\,(\ref{Abb_Datenauswertung_VerteilungDistanzenApr2026})  
bestätigt, in der die Häufigkeits-Verteilung der 
paarweisen Distanzen $R$ dargestellt ist und die ein 
Maximum der euklidischen Distanzen bei rd.\,$7\,km$ 
aufweist. Des Weiteren fällt auf, dass die Anzahl 
der Distanz-Paare für Werte größer als dieses 
Häufigkeits-Maximums bis hin zum maximalen Wert bei etwa 
$R_{max}=22\,km$ sehr schnell abfällt. Daher müssen die \textit{absoluten Häufigkeiten}, 
die bei den nachfolgenden Schätzungen der Semivarianz
einfließen, insbesondere für diese Distanz-Klassen 
genau betrachtet 
werden. Die in der Abbildung ausgewiesene Anzahl $N$ 
der in der Zählung berücksichtigten Paare 
errechnet sich gemäß $N = n\cdot (n-1)/2$, also 
$83\cdot(83-1)/2=3403$ 
(vgl.\,Kapitel \ref{KapitelTheorieVariogrammSchätzung}). \\

Abb.\,(\ref{Abb_Datenauswertung_VerteilungP25MesswerteApr2026}) 
verdeutlicht, dass -wie bei den Einzelmessungen 
(vgl.\,Abb.\,(\ref{Abb_Datenstruktur_Häufigkeitsanalyse_P1_Sensor6117}))- 
auch für die arithmetisch gemittelten Messwerte eine 
Normalverteilungs-Annahme ungerechtfertigt erscheint.

\begin{figure}[H] 
\centering
\includegraphics[scale = 0.45]{Datenauswertung_VerteilungP25MesswerteApr2026} 
\caption{\small{Häufigkeits-Verteilung der $\mathrm{P}_{2.5}$ Messwerte 
im April 2026}}
\label{Abb_Datenauswertung_VerteilungP25MesswerteApr2026}
\end{figure} 










